% !TEX program = xelatex
% ============================================
% 李久阳 · 个人简历（最小可编译版）
% ============================================
% 【重要】Overleaf 操作步骤：
% 1. 全选并删除 Overleaf 左侧编辑器里所有旧代码
% 2. 把本文件完整粘贴进去
% 3. 点左上角 Menu → Compiler 下拉 → 选 XeLaTeX
% 4. 点 Recompile
% 5. 如果还报错，把报错原样贴给我
% ============================================

\documentclass[12pt,a4paper]{ctexart}

% ---- 页面 ----
\usepackage[top=2cm, bottom=2cm, left=2.2cm, right=2.2cm]{geometry}

% ---- 颜色 ----
\usepackage{xcolor}
\definecolor{primary}{HTML}{2c3e50}
\definecolor{accent}{HTML}{2980b9}
\definecolor{gray}{HTML}{7f8c8d}

% ---- 超链接 ----
\usepackage[hidelinks]{hyperref}
\hypersetup{colorlinks=true, urlcolor=accent}

% ---- 自定义命令 ----
\newcommand{\cvsection}[1]{
  \vspace{0.8em}
  \noindent
  {\color{primary}\large\textbf{#1}}
  \vspace{0.15em}
  \hrule
  \vspace{0.5em}
}

\begin{document}

\begin{center}
  {\color{primary}\huge\textbf{李久阳}} \\[0.3em]
  {\large 应用物理学本科生} \\[0.4em]
  {\color{gray}\small jiuyang@example.com \quad github.com/jiuyang \quad 中国·北京}
\end{center}

\cvsection{教育背景}
\textbf{XX大学 · 应用物理学} \hfill 2023--2027（预计）

主修：量子力学、统计物理、电动力学、固体物理、数学物理方法

自学：信息论、认知科学、决策论、复杂系统

\cvsection{研究兴趣}
\begin{itemize}
  \item 信息论与社会系统：将 Shannon 框架推广至注意力分配与认知熵
  \item 认知协作系统：LLM 时代下人-AI 协作的架构设计
  \item 复杂系统建模：从物理学出发，为社会学、认知科学建立形式模型
\end{itemize}

\cvsection{项目经验}
\textbf{笔默 —— AI for Reading} \hfill 2026.03--至今

纸质阅读的"第二层智能"——AI 贴附在阅读行为链路上，不替代，只增强。独立完成商业计划书（十三章），设计硬件形态、AI 交互范式与社区生态。

\textbf{忆阻器时间展宽微弱电流测量} \hfill 2026

利用忆阻器非易失性电荷积分特性，将 pA 级微弱电流测量转化为阻值/时间测量。

\cvsection{技能}
系统化分析 · 跨域建模 · 协议设计 · 信息论 · 认知科学 · 边界条件穷举

Python · Markdown · Git · LaTeX(学习中) · 物理实验设计 · 数学推导

\cvsection{方法论工具箱}
\textbf{分析}：公理化驱动 · 边界条件穷举 · 比较矩阵 · 元层次追问

\textbf{建构}：协议化思维 · 迭代式精炼 · 递归式自我应用 · 反玄学五步法

\cvsection{语言}
中文（母语）\hfill 英文（学术阅读与写作）\hfill 德语（入门）

\vspace{1em}
\begin{center}
  {\color{gray}\footnotesize 「我相信，人类的好奇心可以延向多远的地方。」}
\end{center}

\end{document}
